\documentclass[../main]{subfiles}
\begin{document}
\chapter{Entropía}

\img{images/09/entropia.jpg}{¿La entropía es desorden?}

\section{Entropía}

La \textbf{entropía} es una propiedad termodinámica muy útil . Ha sido definida por 

\info{\emph{Clausius} de la siguiente manera}{
\emph{Si en la transformación reversible una pequeña cantidad de calor absorbido se la divide por la temperatura absoluta a la cual se la toma, el resultado es igual a la correspondiente variación de entropía}
}

\begin{equation}
    \Delta S  =\dfrac{\Delta Q}{T} 
\end{equation}


\actividad{ Investigar y Responder}
\begin{enumerate}
    \item ¿Qué es la muerte térmica del universo?
    \item Entropía. Definir.
    \item ¿Quién fue \emph{Rudolf Clausius} y qué enunció?
    \item ¿Qué es un proceso térmico reversible e irreversible? ¿Porqué?
    \item ¿Qué es un proceso isotérmico?
    \item ¿Qué enuncia el segundo principio de la termodinámica?
    \item ¿Qué es la entropía de la información y que relación tiene con la entropía térmica?
\end{enumerate}

\actividad{ Resolver los siguientes problemas de entropía.}
\begin{enumerate}
    \item El aire de una máquina térmica que evoluciona con un ciclo ideal de CARNOT ingresa a $327^\circ C$ y escapa a $27^\circ C$.
    Suponiendo que la cantidad de calor entregada al aire es de $60.000kcal$, calcular:
    \begin{enumerate}
        \item La variación de entropía que experimetó el aire.
        \item La cantidad de calor cedida al foco frío.
        \item El rendimiento térmico de la máquina.
    \end{enumerate}
    \item Calcular la variación de entropía que experimenta una masa de aire de $100kg$ de peso, encerrada en un recipiente rígido, cuando se calienta desde $27^\circ C$ a $177^\circ C$.
    \item Calcular la variación de entropía que experimenta una masa de helio de 10kg de peso, encerrada en un globo a la presión de 1 atmósfera, cuando se enfría desde $57^\circ C$ a $7^\circ$.
    \item $100 m^3$ de Hidrógeno se expanden isotérmicamente, a la temperatura de $87^\circ C$ desde la presión de 5 atmósferas a la presiónd e 1 atmósfera, calcular:
    \begin{enumerate}
        \item la variación de entropía que experimenta la masa de hidrógeno.
        \item El volumen final.
        \item La cantidad de calor entregada.
    \end{enumerate}
    \item Haciendo uso de la tabla 10 de libro, calcular la variación de \emph{entropía} que experimenta una masa de aire de $1kg$ de peso cuando se calienta a presión constante desde $60^\circ C$ hasta $282^\circ C$.
    \item Durante un proceso de adición de calor isotérmico de un ciclo de Carnot, se agregan 900kJ al fluido de trabajo de una fuente que está a $\SI{400}{\celsius}$. Determinar:
    \begin{enumerate}
        \item Cambio de entropía del fluido.
        \item Cambio de entropía de la fuente.
        \item Cambio de entropía para el total del proceso.
    \end{enumerate}
    \item El gas refrigerante R134A entra como mezcla saturada de líquido vapor a una presión de $160kPa$ en el serpentín de un evaporador de un sistema de refrigeración. El refrigerante absorbe $180kJ$ de calor del espacio refrigerado que mantiene $\SI{-5}{\celsius}$ y sale como vapor saturado a la misma presión. Determinar:
    \begin{enumerate}
        \item Cambio de entropía del refrigerante.
        \item cambio de entropía del espacio refrigerado.
        \item cambio de entropía total.
    \end{enumerate}
    
    \item Un recipiente rígido contiene inicialmente $5kg$ de refrigerante desconocido, a $\SI{20}{\celsius}$ y $140kpa$, la sustancia se enfría mientras es agitado hasta que su presión disminuye a $100kPa$. Determine el cambio de entropía del refrigerante luego de comprimirse, se licua calor latente de $L_v=217 kJ/kg$.
    
    
    \item Un recipiente rígido contiene inicialmente $4kg$ de refrigerante R134A, a $\SI{6}{\celsius}$ y $140kpa$, la sustancia se enfría mientras es agitado hasta que su presión disminuye a $100kPa$. Determine el cambio de entropía del refrigerante durante el proceso.
    
    
    \item Un recipiente rígido contiene inicialmente $5kg$ de refrigerante R134A, a $\SI{20}{\celsius}$ y $140kpa$, la sustancia se enfría mientras es agitado hasta que su presión disminuye a $100kPa$. Determine el cambio de entropía del refrigerante durante el proceso.
    
    
    
    \item Un dispositivo rígido de $ 1m^3$ de volumen contiene agua a 30bar y $\SI{600}{\celsius}$. El agua se enfría y la presión baja a 15 bar, el ambiente que recibe el calor está a 1.1bar y $\SI{27}{\celsius}$ . Determine:
    \begin{enumerate}
        \item Variación de la entropía del agua.
        \item Producción total de entropía en el proceso.
        \item dibujar el proceso en un diagrama T-S.
    \end{enumerate}
\end{enumerate}


\actividad{Diagrama de Mollier}

%    \begin{table}[]
 %       \centering
  %  \scriptsize
  %  \begin{tabular}{|l|c|c|c|c|c|c|c|} \hline
   %  $T (^\circ C$) & P(kPa) & $\rho_l (kg/m^3)$ & $V_g (m^3/kg)$ & $h_l (kJ/kg)$ & $h_g (kJ/kg)$ & $S_l(kJ/kg ^\circ K)$ &$S_g(kJ/kg ^\circ K)$
     
    %\csvreader[head to column names]{tabla.csv}{}
    %{\\ \hline \temperatura\ & \presion & %\densidadl&\volumeng&\entalpial&\entalpiag&\entropial&\entropiag}\\ \hline
    
   % \end{tabular}
    %    \caption{ Propiedades del líquido y vapor saturados para el R-134A}
   % \end{table}
    
    
    
\end{document}

\end{document}